\begin{definition}[Region two-block tensor]%
    \label{def:peps_regionTwoBlock}
    \lean{TNLean.PEPS.regionTwoBlock}
    \leanok
    For a finite region $R\subseteq V$, the blocked tensor of $R$ may be read
    as a two-block tensor with a one-point external boundary and shared bonds
    the edges crossing the boundary of $R$:
    \[
        B_{A,R}(*,\eta,\sigma)=T_{A,R}^{\eta}(\sigma).
    \]
    This is the region block used in the one-block-against-complement step of
    \cite[Section~3]{Molnar2018NormalPEPS}.
\end{definition}

\begin{definition}[Complement-boundary identification]%
    \label{def:peps_regionComplementBoundaryConfig}
    \lean{TNLean.PEPS.regionBoundaryEdgeToCompl,
        TNLean.PEPS.regionBoundaryEdgeComplEquiv,
        TNLean.PEPS.regionComplementBoundaryConfig}
    \leanok
    The edges crossing the boundary of $R$ are the same edges crossing the
    boundary of $V\setminus R$. Thus a boundary configuration $\eta$ for $R$
    determines the boundary configuration $\eta^c$ for the complement by reading
    the same edge labels through this identification.
\end{definition}

\begin{definition}[Complement-region two-block tensor]%
    \label{def:peps_regionComplementTwoBlock}
    \lean{TNLean.PEPS.regionComplementTwoBlock}
    \leanok
    \uses{def:peps_regionComplementBoundaryConfig}
    For a finite region $R\subseteq V$, the complement block $V\setminus R$ is
    a two-block tensor over the same crossing bonds:
    \[
        B_{A,R^c}(*,\eta,\tau)
        =
        T_{A,V\setminus R}^{\eta^c}(\tau).
    \]
    This is the complementary block in the region/complement comparison of
    \cite[Section~3]{Molnar2018NormalPEPS}.
\end{definition}

\begin{theorem}[Injectivity of the complement-region two-block tensor]%
    \label{thm:peps_isTwoBlockInjective_regionComplementTwoBlock}
    \lean{TNLean.PEPS.isTwoBlockInjective_regionComplementTwoBlock}
    \leanok
    \uses{def:peps_regionComplementTwoBlock, def:peps_regionInjectivityData}
    If the blocked tensor family of $V\setminus R$ is injective, then the
    complement-region two-block tensor is injective:
    \[
        \operatorname{inj}(T_{A,V\setminus R})
        \Longrightarrow
        \operatorname{inj}(B_{A,R^c}).
    \]
\end{theorem}
\begin{proof}\leanok
    The one-point external boundary does not change the coefficient family, and
    the complement-boundary identification only reindexes the crossing bonds.
    More explicitly, if $\eta_1,\eta_2$ are boundary configurations on
    $\partial R$, then
    \[
        \eta_1^c=\eta_2^c\quad\Longrightarrow\quad \eta_1=\eta_2 .
    \]
    Thus the coefficients of $T_{A,V\setminus R}$ and $B_{A,R^c}$ differ only
    by an injective relabelling of boundary configurations.
\end{proof}

\begin{definition}[Region insertion coefficient]%
    \label{def:peps_regionInsertedCoeff}
    \lean{TNLean.PEPS.regionInsertedCoeff}
    \leanok
    \uses{def:peps_regionTwoBlock, def:peps_regionComplementTwoBlock}
    Let $f$ be an edge crossing the boundary of $R$, and let
    $M\in\MN{D_f}$.  The region insertion coefficient is
    \[
        C_A(R,f,M;\sigma,\tau)
        =
        \sum_{\substack{\mu,\nu\\
        \mu|_{\partial R\setminus\{f\}}
        =\nu|_{\partial R\setminus\{f\}}}}
        M_{\mu_f,\nu_f}\,
        T_{A,R}^{\mu}(\sigma)\,
        T_{A,V\setminus R}^{\nu^c}(\tau).
    \]
    It is the coefficient obtained by inserting $M$ on the boundary bond $f$
    between the region $R$ and its complement.
\end{definition}

\begin{theorem}[Region insertion coefficients determine the inserted matrix]%
    \label{thm:peps_regionInsertedCoeff_injective}
    \lean{TNLean.PEPS.regionInsertedCoeff_injective}
    \leanok
    \uses{def:peps_regionInsertedCoeff}
    Suppose that the blocked tensor maps of $A$ on $R$ and on
    $V\setminus R$ are injective, and that every bond dimension of $A$ is
    positive.  If $f\in\partial R$ and
    $M,M'\in\MN{D_A(f)}$ satisfy
    \[
        C_A(R,f,M;\sigma,\tau)
        =
        C_A(R,f,M';\sigma,\tau)
    \]
    for every physical configuration $\sigma$ on $R$ and $\tau$ on
    $V\setminus R$, then $M=M'$.
\end{theorem}
\begin{proof}
    \leanok
    Subtracting the two coefficients gives
    \[
        C_A(R,f,M-M';\sigma,\tau)=0
    \]
    for every $\sigma,\tau$.  Injectivity of the region and complement
    blocked tensor maps reads off every entry of $M-M'$ from configurations
    agreeing away from $f$, so $M-M'=0$ and hence $M=M'$.
\end{proof}

\begin{definition}[Bond-inserted region insert]%
    \label{def:peps_bond_inserted_region_insert}
    \lean{TNLean.PEPS.bondInsertedRegionInsert}
    \leanok
    Let $f$ be an edge crossing the boundary of $R$, and let
    $M\in\MN{D_f}$.  The bond-inserted region insert
    $I_{A,R,f,M}$ is the insert on $R$ defined by
    \[
        (I_{A,R,f,M})_{\nu,\sigma}
        =
        \sum_{\substack{\mu\in\partial R\\
        \mu|_{\partial R\setminus\{f\}}
        =\nu|_{\partial R\setminus\{f\}}}}
        M_{\mu_f,\nu_f}\,T_{A,R}^{\mu}(\sigma).
    \]
    Thus the complement boundary value $\nu$ is fixed away from $f$, while
    the matrix $M$ couples the two $f$-labels.  This is the one-bond insertion
    in the region/complement coefficient of
    \cite[Section~3]{Molnar2018NormalPEPS}.
\end{definition}

\begin{theorem}[Bond insertion as a deformed region state]
    \label{thm:peps_deformedRegionState_bondInsertedRegionInsert}
    \lean{TNLean.PEPS.deformedRegionState_bondInsertedRegionInsert}
    \leanok
    \uses{def:peps_bond_inserted_region_insert,
        def:peps_regionInsertedCoeff}
    For every physical configuration $\sigma$ on $R$ and $\tau$ on
    $V\setminus R$, the deformed state of the bond-inserted region insert is
    the region insertion coefficient:
    \[
        \Psi_{A,R,I_{A,R,f,M}}(\sigma,\tau)
        =
        C_A(R,f,M;\sigma,\tau).
    \]
\end{theorem}
\begin{proof}
    \leanok
    Expanding the deformed state gives
    \[
        \Psi_{A,R,I_{A,R,f,M}}(\sigma,\tau)
        =
        \sum_{\nu\in\partial R}
        I_{A,R,f,M}(\nu,\sigma)\,
        T_{A,V\setminus R}^{\nu^c}(\tau).
    \]
    Substituting the definition of $I_{A,R,f,M}$ and exchanging the two finite
    sums gives
    \[
    \begin{aligned}
        \Psi_{A,R,I_{A,R,f,M}}(\sigma,\tau)
        &=
        \sum_{\nu\in\partial R}
        \left(
        \sum_{\substack{\mu\in\partial R\\
        \mu|_{\partial R\setminus\{f\}}
        =\nu|_{\partial R\setminus\{f\}}}}
        M_{\mu_f,\nu_f}\,T_{A,R}^{\mu}(\sigma)
        \right)
        T_{A,V\setminus R}^{\nu^c}(\tau)\\
        &=
        C_A(R,f,M;\sigma,\tau).
    \end{aligned}
    \]
\end{proof}

For a region $R$, a physical configuration $\sigma$ on $R$, and a physical
configuration $\tau$ on $V\setminus R$, write $\omega_{R,\sigma,\tau}$ for the
physical configuration on $V$ with
\[
    \omega_{R,\sigma,\tau}(v)=
    \begin{cases}
        \sigma(v),& v\in R,\\
        \tau(v),& v\notin R.
    \end{cases}
\]

\begin{theorem}[Region insertion as an edge insertion]
    \label{thm:peps_regionInsertedCoeff_eq_smul_edgeInsertedCoeff}
    \lean{TNLean.PEPS.regionInsertedCoeff_eq_smul_edgeInsertedCoeff}
    \leanok
    \uses{def:peps_regionInsertedCoeff, def:peps_edgeInsertedCoeff}
    Let $f$ be a boundary edge of $R$, let $N\in\MN{D_B(f)}$, let $\sigma$
    be a physical configuration on $R$, and let $\tau$ be a physical
    configuration on $V\setminus R$.  Put
    \[
        p_B(R)=\prod_{g\notin\partial R}D_B(g).
    \]
    Then the region insertion coefficient factors through the edge-inserted
    coefficient:
    \[
        C_B(R,f,N;\sigma,\tau)
        =
        p_B(R)\,
        E_B\bigl(f,\omega_{R,\sigma,\tau},\mathcal O_{R,f}(N)\bigr).
    \]
    Here $\mathcal O_{R,f}$ is the boundary-edge orientation, equal to the
    identity if the left endpoint of $f$ lies in $R$ and to transpose otherwise.
\end{theorem}
\begin{proof}
    \leanok
    Let $\mathcal C_{R,f}(\sigma,\tau)$ denote the open-edge configurations
    compatible with the two physical configurations and with the two boundary
    readings agreeing away from $f$.  For
    $\xi\in\mathcal C_{R,f}(\sigma,\tau)$, let
    $\operatorname{Fib}_{R,f}(\xi)$ be the corresponding fibre of agreeing
    boundary readings.  Grouping the expansion of $C_B(R,f,N;\sigma,\tau)$ by
    such an open-edge configuration $\xi$ gives
    \[
    \begin{aligned}
        C_B(R,f,N;\sigma,\tau)
        &=
        \sum_{\xi\in\mathcal C_{R,f}(\sigma,\tau)}
        \#\operatorname{Fib}_{R,f}(\xi)\,
        N_{\xi_{R,f}^{(1)},\xi_{R,f}^{(2)}}
        \prod_{v\in V}
        B_v\bigl(\xi_v,\omega_{R,\sigma,\tau}(v)\bigr),\\
        \#\operatorname{Fib}_{R,f}(\xi)
        &=
        \prod_{g\notin\partial R}D_B(g)
        =
        p_B(R).
    \end{aligned}
    \]
    Substituting this fibre cardinality gives
    \[
    \begin{aligned}
        C_B(R,f,N;\sigma,\tau)
        &=
        p_B(R)
        \sum_{\xi\in\mathcal C_{R,f}(\sigma,\tau)}
        N_{\xi_{R,f}^{(1)},\xi_{R,f}^{(2)}}
        \prod_{v\in V}
        B_v\bigl(\xi_v,\omega_{R,\sigma,\tau}(v)\bigr)\\
        &=
        p_B(R)\,
        E_B\bigl(f,\omega_{R,\sigma,\tau},\mathcal O_{R,f}(N)\bigr).
    \end{aligned}
    \]
    The second equality is the definition of the edge-inserted coefficient,
    with the inserted matrix read in the boundary-edge orientation.
\end{proof}

\begin{theorem}[Bond insertion through a corner extension]
    \label{thm:peps_regionInsertedCoeff_eq_extendInsert_bondInserted}
    \lean{TNLean.PEPS.regionInsertedCoeff_eq_extendInsert_bondInserted}
    \leanok
    \uses{thm:peps_deformedRegionState_bondInsertedRegionInsert,
        thm:peps_corner_extension_state_preserves}
    Let $R\subseteq S$, and suppose that all bond dimensions are positive.
    For every global physical configuration $\sigma$,
    \[
        C_A\bigl(R,f,M;\sigma|_R,\sigma|_{V\setminus R}\bigr)
        =
        \Psi_{A,S,
        \operatorname{Ext}_{R\subseteq S}(I_{A,R,f,M})}(\sigma).
    \]
    In particular, the one-bond insertion on $R$ may be read after extending
    the corresponding insert to any larger region $S$.
\end{theorem}
\begin{proof}
    \leanok
    The preceding theorem and state preservation for corner extension give
    \[
    \begin{aligned}
        C_A\bigl(R,f,M;\sigma|_R,\sigma|_{V\setminus R}\bigr)
        &=
        \Psi_{A,R,I_{A,R,f,M}}
            (\sigma|_R,\sigma|_{V\setminus R})\\
        &=
        \Psi_{A,S,
            \operatorname{Ext}_{R\subseteq S}(I_{A,R,f,M})}(\sigma).
    \end{aligned}
    \]
\end{proof}

\begin{theorem}[Region insertion as a two-block insertion]%
    \label{thm:peps_twoBlockInsertedCoeff_eq_regionInsertedCoeff}
    \lean{TNLean.PEPS.twoBlockInsertedCoeff_eq_regionInsertedCoeff}
    \leanok
    \uses{def:peps_regionInsertedCoeff, def:peps_regionTwoBlock,
        def:peps_regionComplementTwoBlock}
    The abstract two-block insertion coefficient of the pair
    $(B_{A,R},B_{A,R^c})$ is the region insertion coefficient:
    \[
        C_{B_{A,R},B_{A,R^c}}(f,M;*,*,\sigma,\tau)
        =
        C_A(R,f,M;\sigma,\tau).
    \]
\end{theorem}
\begin{proof}\leanok
    Expanding the abstract two-block coefficient gives exactly the two boundary
    sums in the displayed formula for $C_A(R,f,M;\sigma,\tau)$.
\end{proof}

\begin{theorem}[Region insertion under bond-dimension transport]%
    \label{thm:peps_regionInsertedCoeff_reindexTensor}
    \lean{TNLean.PEPS.regionInsertedCoeff_reindexTensor}
    \leanok
    \uses{def:peps_regionInsertedCoeff}
    If a tensor $B$ is transported along an equality of bond dimensions, then
    the region insertion coefficient is transported by the induced
    matrix-algebra identification on the inserted bond:
    \[
        C_{B^h}(R,f,N;\sigma,\tau)
        =
        C_B(R,f,\phi_f(N);\sigma,\tau).
    \]
\end{theorem}
\begin{proof}\leanok
    Reindex both boundary-configuration sums by the bond-dimension
    equivalences. The constraint away from $f$ and the two blocked-region
    weights are preserved, while the inserted matrix entry becomes the
    transported entry $\phi_f(N)$.
\end{proof}

\begin{theorem}[Same two-block insertions from region insertions]%
    \label{thm:peps_sameTwoBlockInsertions_of_regionInsertedCoeff_eq}
    \lean{TNLean.PEPS.sameTwoBlockInsertions_of_regionInsertedCoeff_eq}
    \leanok
    \uses{thm:peps_twoBlockInsertedCoeff_eq_regionInsertedCoeff,
        thm:peps_regionInsertedCoeff_reindexTensor}
    Suppose $A$ and $B$ have identified bond dimensions. Write $B^h$ for the
    transported tensor, and write $\phi_f$ for the induced transport on the
    boundary bond $f$. If
    \[
        C_A(R,f,N;\sigma,\tau)
        =
        C_B(R,f,\phi_f(N);\sigma,\tau)
    \]
    for every boundary edge $f$, matrix $N$, and physical configurations
    $\sigma,\tau$, then the two pairs
    \[
        (B_{A,R},B_{A,R^c}),\qquad
        (B_{B^h,R},B_{B^h,R^c})
    \]
    have the same one-bond insertion coefficients after the same
    bond-dimension transport.
\end{theorem}
\begin{proof}\leanok
    Rewrite the two abstract insertion coefficients as region insertion
    coefficients and use the bond-dimension transport formula:
    \[
        C_{B_{A,R},B_{A,R^c}}(f,N;*,*,\sigma,\tau)
        = C_A(R,f,N;\sigma,\tau)
        = C_B(R,f,\phi_f(N);\sigma,\tau)
        = C_{B_{B^h,R},B_{B^h,R^c}}(f,N;*,*,\sigma,\tau).
    \]
\end{proof}

\begin{theorem}[Swapping the two blocks in a one-bond insertion]
    \label{thm:peps_twoBlockInsertedCoeff_swap}
    \lean{TNLean.PEPS.twoBlockInsertedCoeff_swap}
    \leanok
    Let \(B_1,B_2\) be two blocks joined along a shared bond \(f\).  Then inserting
    a matrix \(X\) while contracting \(B_1\) against \(B_2\) is the same as
    contracting \(B_2\) against \(B_1\) with the transposed insertion:
    \[
        C_{B_1,B_2}(f,X;\eta_1,\eta_2,\sigma_1,\sigma_2)
        =
        C_{B_2,B_1}(f,X^T;\eta_2,\eta_1,\sigma_2,\sigma_1).
    \]
\end{theorem}
\begin{proof}\leanok
    Exchange the two boundary-configuration sums.  The condition that the two
    configurations agree away from \(f\) is symmetric, and
    \(X^T_{\nu_f,\mu_f}=X_{\mu_f,\nu_f}\).
\end{proof}

\begin{theorem}[Complement-side reading of a region insertion]
    \label{thm:peps_regionInsertedCoeff_complement_reading}
    \lean{TNLean.PEPS.regionInsertedCoeff_eq_complementTwoBlock}
    \leanok
    \uses{thm:peps_twoBlockInsertedCoeff_swap,
        thm:peps_twoBlockInsertedCoeff_eq_regionInsertedCoeff}
    For a boundary edge \(f\in\partial R\), the region-inserted coefficient can be
    read by contracting the complement block first:
    \[
        C_A(R,f,X;\sigma,\tau)
        =
        C_{A,R^c;R}(f,X^T;\tau,\sigma).
    \]
    Here \(C_{A,R^c;R}\) denotes the same one-bond insertion coefficient with the
    two blocks \(R^c\) and \(R\) interchanged.
\end{theorem}
\begin{proof}\leanok
    Identify \(C_A(R,f,X;\sigma,\tau)\) with the two-block insertion of
    \((B_{A,R},B_{A,R^c})\), then apply
    Theorem~\ref{thm:peps_twoBlockInsertedCoeff_swap}.
\end{proof}

\begin{theorem}[Incident form gives the transferred endpoint form]
    \label{thm:peps_hform_of_region_incident_form}
    \lean{TNLean.PEPS.hform_of_isIncidentMatrixForm}
    \leanok
    \uses{thm:peps_physical_to_virtual_insertion}
    Assume that \(A\)'s and \(B\)'s components at the in-region endpoint of \(f\)
    are linearly independent, and that \(D_B(e)>0\) for every edge \(e\).
    Let \(W\) be the virtual pullback, at the in-region endpoint of \(f\), of the
    physical operator obtained by inserting \(X^T\) on \(f\).  If
    \[
        W=\mathcal I_f(P)
    \]
    is in incident-matrix form on the boundary leg \(f\), then the transfer matrix
    read off from \(W\) reinserts to \(W\).
\end{theorem}
\begin{proof}\leanok
    The transfer matrix is defined by reading the incident matrix out of \(W\).
    Reading \(\mathcal I_f(P)\) gives \(P\), and reinserting \(P\) returns \(W\).
\end{proof}

\begin{theorem}[Virtual pullback realizes the endpoint operator]
    \label{thm:peps_region_insertion_virtual_pullback_realizes}
    \lean{TNLean.PEPS.localTensorMap_localVirtualOpOfPhysicalOpAt_regionInsertionOp}
    \leanok
    \uses{def:peps_same_state, def:IsVertexInjective, def:peps_localTensorMap,
        thm:peps_physical_to_virtual_insertion}
    Assume that \(A\) and \(B\) have the same PEPS state, are vertex-injective,
    and have positive bond dimensions.  Let \(v\) be the in-region endpoint of
    \(f\), and let \(W\) be the virtual pullback, for \(B\) at \(v\), of the
    physical operator obtained by inserting \(X^T\) on \(f\).  Then for every
    local virtual coefficient \(c\),
    \[
        \mathcal T_{B,v}(Wc)
        =
        O_{A,R,f}(X^T)\,\mathcal T_{B,v}(c).
    \]
\end{theorem}
\begin{proof}\leanok
    The physical endpoint operator preserves the image of
    \(\mathcal T_{B,v}\). The chosen left inverse defining the virtual pullback
    therefore realizes the same action after applying \(\mathcal T_{B,v}\).
\end{proof}

\begin{definition}[Out-of-region endpoint and blocked-region inverse]
    \label{def:peps_region_out_endpoint_block_inverse}
    \lean{TNLean.PEPS.regionBoundaryEdgeOutVertex,
        TNLean.PEPS.regionBoundaryEdgeOutVertex_not_mem,
        TNLean.PEPS.regionBoundaryEdgeOutVertex_mem_compl,
        TNLean.PEPS.regionBoundaryEdgeInVertex_compl_eq_outVertex,
        TNLean.PEPS.regionBlockedTensorMap,
        TNLean.PEPS.regionBlockedTensorMap_injective_of_injective,
        TNLean.PEPS.regionBlockedTensorMap_ker_eq_bot_of_injective,
        TNLean.PEPS.regionBlockedLeftInverse,
        TNLean.PEPS.regionBlockedLeftInverse_comp_regionBlockedTensorMap,
        TNLean.PEPS.regionBlockedLeftInverse_apply_regionBlockedTensorMap,
        TNLean.PEPS.regionBlockedTensorMap_apply,
        TNLean.PEPS.regionBlockedTensorMap_single,
        TNLean.PEPS.regionBlockedLeftInverse_regionBlockedWeight}
    \leanok
    \uses{def:peps_regionInjectivityData}
    The out-of-region endpoint of \(f\in\partial R\) is the endpoint outside
    \(R\).  It is the in-region endpoint of the same edge viewed as an edge of
    \(\partial(V\setminus R)\).  The blocked-region map is
    \[
        c\longmapsto \sum_{\mu\in\partial R} c_\mu\,T_{A,R}^{\mu},
    \]
    and injectivity of the blocked tensor gives a left inverse.  In particular,
    \[
        L_{A,R}\!\left(T_{A,R}^{\mu}\right)=e_\mu .
    \]
\end{definition}

\begin{theorem}[Complement-side cast and endpoint realization]
    \label{thm:peps_regionInsertedCoeff_complement_cast}
    \lean{TNLean.PEPS.mem_compl_compl_iff,
        TNLean.PEPS.regionDoubleComplVertexEquiv,
        TNLean.PEPS.regionDoubleComplPhysicalConfig,
        TNLean.PEPS.isRegionBoundaryEdge_compl_compl_iff,
        TNLean.PEPS.regionBoundaryEdgeDoubleComplEquiv,
        TNLean.PEPS.regionDoubleComplBoundaryConfig,
        TNLean.PEPS.regionBlockedWeight_doubleCompl,
        TNLean.PEPS.regionComplementBoundaryConfigEquiv,
        TNLean.PEPS.regionComplementBoundaryConfig_apply_toCompl,
        TNLean.PEPS.regionComplementBoundaryConfig_compl_compl,
        TNLean.PEPS.sameAwayFromBond_complementBoundaryConfig,
        TNLean.PEPS.regionInsertedCoeff_eq_compl,
        TNLean.PEPS.regionInsertedCoeff_eq_regionRealizationSum_compl,
        TNLean.PEPS.regionInsertedCoeff_eq_smul_op_regionStateVec_compl}
    \leanok
    \uses{def:peps_region_out_endpoint_block_inverse,
        thm:peps_regionInsertedCoeff_complement_reading}
    With \(f'\) the boundary edge \(f\) read on \(V\setminus R\), one has
    \[
        C_A(R,f,X;\sigma,\tau)
        =
        C_A(V\setminus R,f',X^T;\tau,\sigma^{cc}),
    \]
    where \(\sigma^{cc}\) is \(\sigma\) transported across
    \(V\setminus(V\setminus R)=R\).  Consequently the same coefficient is realized
    through the endpoint of \(f\) lying outside \(R\).
\end{theorem}
\begin{proof}\leanok
    Reindex the two boundary-configuration sums by the boundary-edge equivalence
    between \(R\) and its complement.  The double complement carries the blocked
    weight of \(V\setminus(V\setminus R)\) back to the blocked weight of \(R\).
    Applying the usual endpoint realization to \(V\setminus R\) gives the displayed
    out-of-region endpoint realization.
\end{proof}

\begin{theorem}[Region resonate identity and blocked read-offs]
    \label{thm:peps_region_resonate_identity_readoffs}
    \lean{TNLean.PEPS.regionInsertionOp_regionStateVec_pin_compl,
        TNLean.PEPS.region_resonate_identity,
        TNLean.PEPS.regionComplementRow,
        TNLean.PEPS.regionInsertedCoeff_eq_complement_blockedMap,
        TNLean.PEPS.regionRegionRow,
        TNLean.PEPS.regionInsertedCoeff_eq_region_blockedMap,
        TNLean.PEPS.regionBlockedLeftInverse_complement_regionInsertedCoeff,
        TNLean.PEPS.regionBlockedLeftInverse_region_regionInsertedCoeff}
    \leanok
    \uses{def:peps_region_out_endpoint_block_inverse,
        thm:peps_regionInsertedCoeff_complement_cast}
    Assume that \(A\)'s components at the in-region endpoint of \(f\) and at the
    in-complement-region endpoint of the complementary boundary edge are linearly
    independent, and that \(A\) and \(B\) define the same PEPS state.  Let
    \(P_R\) and \(P_{R^c}\) be the interior bond products of \(R\) and \(R^c\).
    Let \(O^{\mathrm{in}}_f(X)\) and \(O^{\mathrm{out}}_f(X)\) be the physical
    endpoint operators obtained from the two virtual insertions, and let
    \(\Psi_B(R;\sigma,\tau)\) and \(\Psi_B(R^c;\tau,\sigma^{cc})\) be the
    corresponding closed \(B\)-state vectors.  Then the two endpoint readings of
    the same inserted coefficient give
    \[
        P_R\,O^{\mathrm{in}}_f(X)\,\Psi_B(R;\sigma,\tau)
        =
        P_{R^c}\,O^{\mathrm{out}}_f(X)\,\Psi_B(R^c;\tau,\sigma^{cc}).
    \]
    With \(\rho^{R^c}_{X,\sigma}\) and \(\rho^{R}_{X,\tau}\) denoting the two
    row functions, one has
    \[
        C_A(R,f,X;\sigma,-)
        =
        T_{A,R^c}\!\left(\rho^{R^c}_{X,\sigma}\right),
        \qquad
        C_A(R,f,X;-,\tau)
        =
        T_{A,R}\!\left(\rho^{R}_{X,\tau}\right).
    \]
    Hence the blocked left inverses read off the rows:
    \[
        L_{A,R^c}\!\left(C_A(R,f,X;\sigma,-)\right)
        =
        \rho^{R^c}_{X,\sigma},
        \qquad
        L_{A,R}\!\left(C_A(R,f,X;-,\tau)\right)
        =
        \rho^{R}_{X,\tau}.
    \]
\end{theorem}
\begin{proof}\leanok
    The in-region and out-of-region endpoint pins both evaluate to
    \(C_A(R,f,X;\sigma,\tau)\), hence they agree.  Expanding
    \(C_A(R,f,X;\sigma,\tau)\) as a sum over boundary configurations isolates,
    respectively, the \(R^c\)-blocked and \(R\)-blocked tensor maps; the left
    inverses then read off the row functions.
\end{proof}

